\documentclass[12pt]{article}
\usepackage[utf8]{inputenc}
\usepackage{geometry}
\geometry{a4paper, margin=1in}
\usepackage{booktabs}
\usepackage{enumitem}
\usepackage{titlesec}
\usepackage{amsmath}
\usepackage{amssymb}
\usepackage{graphicx}
\usepackage{xcolor}
\usepackage{hyperref}

\titleformat{\section}{\large\bfseries}{\thesection.}{0.5em}{}
\titleformat{\subsection}{\normalsize\bfseries}{\thesubsection}{0.5em}{}

% Placeholder marker: delete every \todo{} before submitting.
\newcommand{\todo}[1]{\textcolor{red}{[TODO: #1]}}

\title{\textbf{Preliminary Report: AI Racing Challenge} \\ \large Pure Pursuit vs.\ MPCC vs.\ PPO}
\author{\todo{Nataly, Berinay, Christian, Pavan}}
\date{\today}

\begin{document}
\maketitle

\noindent\textbf{GitHub repository:} \url{https://github.com/wukongxzero/AI-racing-challenge}

%=====================================================================
\section{Summary}
% 3-5 sentences: what the project is, what is working today, what is next.
% [PROMPT: "a short description of the project"] reuse ppo_plan.tex §1.

\todo{summary}

%=====================================================================
\section{Changes Since the Proposal}
% What changed in scope or design since 11-08, and why.
% (e.g. racing objective vs. path tracking, tire model, MPC model choice, roles)

\todo{changes and reasons}

%=====================================================================
\section{Simulation Environment}
% [PROMPT: "tools to be used (simulator, existing code, etc)"] -- the simulator part.
% The prompt also wants the OTHER tools (SB3, CasADi+IPOPT, Pygame, Gymnasium,
% PyTorch, TensorBoard): add them in Interface or a short table (ppo_plan.tex §2).
\subsection{Vehicle Model}
% State, controls, dynamics model, integrator, dt, parameters.

\todo{vehicle model}

\subsection{Track Representation}
% How the track is stored; what the projection function returns.

\todo{track}

\subsection{Interface Shared by All Controllers}
% The step/reset contract every controller codes against; termination; config file.

\todo{interface}

\subsection{Status}
% What is implemented and tested; what is not yet.

\todo{sim status}

%=====================================================================
\section{Controllers}
% [PROMPT: "algorithms to be implemented"]
\subsection{Pure Pursuit (Baseline)}
% Steering law, lookahead, speed profile, status.

\todo{pure pursuit}

\subsection{MPCC}
% Optimal control problem: decision variables, cost, dynamics constraints,
% inequality constraints, horizon, solver, warm start. Status.

\todo{MPC formulation}

\subsection{PPO}
% Observation, action, reward, termination, network, training setup, seeds. Status.

% Copied word for word from ppo_plan.tex §4 (your text). TODO for you: trim to
% "brief" (the prompt says brief), and add Training + Status below.
\begin{enumerate}[leftmargin=1.5em, itemsep=0.35em]
    \item \textbf{State vs.\ Observation ($x_n$ vs.\ $o_n$):} While the underlying simulator state $x_n$ tracks global coordinates $(X, Y, \psi)$, wheel spins, and track progress, PPO receives a 15-dimensional ego-centric observation vector $o_n \in \mathbb{R}^{15}$ that omits global $(X, Y, \psi, V_X, V_Y)$ so the policy cannot memorize map coordinates and instead generalizes across corners and unseen tracks:
    \begin{itemize}[leftmargin=1.5em, itemsep=0.15em]
        \item \textbf{Body-frame motion (3):} Forward speed $v_x$, lateral slip speed $v_y$, and yaw rate $r$.
        \item \textbf{Actuator state (1):} Actual current steering angle $\delta_{\text{curr}}$ (necessary because steering is rate-limited at $3.2~\text{rad/s}$ and can lag behind commanded steering).
        \item \textbf{Track-relative errors (2):} Signed lateral offset $e_c$ from the centerline and heading error $e_\psi \in [-\pi, \pi]$ relative to the road tangent.
        \item \textbf{Previous action (2):} $[\text{throttle}_{n-1}, \text{steering}_{n-1}]$ to encourage smooth control transitions.
        \item \textbf{Curvature lookahead (6):} Upcoming centerline curvatures $\kappa_1, \dots, \kappa_6$ sampled every $0.5~\text{m}$ from $0.5~\text{m}$ to $3.0~\text{m}$ ahead ($\approx 0.5~\text{s}$ preview at $5\text{--}6~\text{m/s}$).
        \item \textbf{Local track half-width (1):} $w_{\text{half}}$ so the policy knows available road width when width varies between straights and hairpins.
    \end{itemize}
    All 15 inputs are standardized online using Stable-Baselines3's \texttt{VecNormalize} (running mean and variance) so velocities, angles, and curvatures sit on a comparable scale.

    \item \textbf{Control / Action Space ($u_n$):} A 2D continuous command $u_n = [\text{throttle}_n, \text{steering}_n] \in [-1, 1]^2$ issued at $20~\text{Hz}$ and held over each $\Delta t = 50~\text{ms}$ control step. Inside \texttt{RacingEnv}, the shared limiter scales steering to $[-0.42, 0.42]~\text{rad}$ and enforces the $\le 3.2~\text{rad/s}$ steering slew rate.

    \item \textbf{Transition Dynamics:} Treated as a model-free Markov transition $x_{n+1} \sim P(\cdot \mid x_n, u_n)$ stepped by \texttt{RacingEnv} over $2\text{--}5~\text{ms}$ physics sub-steps (\texttt{car="numpy"} during initial bring-up, \texttt{car="mujoco"} for full rigid-body contact and tire slip). Unlike MPCC, PPO requires no explicit bicycle model or fitted tire parameters.

    \item \textbf{Stage Cost / Reward ($r_n$):} At each $50~\text{ms}$ step, the reward balances fast forward progress against stability and track boundaries:
    \[
        r_n = w_{\text{prog}} \Delta s_n - w_c e_c^2 - w_\psi e_\psi^2 - w_{\Delta u} \|u_n - u_{n-1}\|^2 - w_{\text{time}} - R_{\text{crash}} \cdot \mathbb{I}(\text{off-track})
    \]
    \begin{itemize}[leftmargin=1.5em, itemsep=0.15em]
        \item \textbf{Centerline progress ($w_{\text{prog}} \Delta s_n$, $w_{\text{prog}} = 1.0$):} Rewards distance gained along the centerline ($\Delta s_n = (s_n - s_{n-1} + L/2) \bmod L - L/2$ to handle start/finish wrap-around), preventing sideways drifting or slow wide arcs.
        \item \textbf{Alignment penalties ($w_c = w_\psi = 0.05$):} Mildly discourages excessive centerline offset and slip angles while remaining small enough to permit clean apex cutting.
        \item \textbf{Action smoothness ($w_{\Delta u} = 0.05$):} Penalizes step-to-step command jumps to prevent $20~\text{Hz}$ steering chatter from breaking tire grip.
        \item \textbf{Step time penalty ($w_{\text{time}} = 0.01$):} Ensures stalling yields negative reward so minimum-time lap completion is the discounted optimum.
        \item \textbf{Off-track penalty ($R_{\text{crash}} = 10.0$):} Applied upon off-track termination; kept moderate so the policy respects track limits without becoming overly timid in corners.
    \end{itemize}

    \item \textbf{Episode Start \& End ($x_0$, Termination vs.\ Truncation):}
    \begin{itemize}[leftmargin=1.5em, itemsep=0.15em]
        \item \textbf{Initial state distribution ($x_0$):} During training, episodes spawn at uniformly random centerline progress $s_0 \in [0, L)$ with small perturbations ($e_c \in [-0.3, 0.3]~\text{m}$, $e_\psi \in [-0.15, 0.15]~\text{rad}$, $v_x \in [0.5, 2.5]~\text{m/s}$) so PPO learns all corners evenly rather than overfitting to Turn~1. During evaluation, episodes start from rest at the official start line ($s_0 = 0$).
        \item \textbf{Terminated (failure, bootstrap $= 0$):} Off-track ($|e_c| > w_{\text{half}} - w_{\text{car}}/2$), reverse heading ($|e_\psi| > \pi/2$), or stall ($< 0.3~\text{m}$ progress over a $2.0~\text{s}$ / 40-step window).
        \item \textbf{Truncated (time limit, critic bootstrap):} Reaching $T_{\max} = 1000$ steps ($50~\text{s}$ at $20~\text{Hz}$), separated from \texttt{terminated} so surviving to the time limit is bootstrapped rather than penalized as a crash.
    \end{itemize}
\end{enumerate}

\paragraph{Training.}
% SB3 PPO from scratch, Adam lr 3e-4 (SB3 default), >= 3 seeds, TensorBoard;
% one sentence on PPO (clip, critic, GAE) + cite Schulman 2017.
\todo{training}

\paragraph{Status.}
% Design done; first run = toy env, then the stand-in car.
\todo{status}

%=====================================================================
\section{Preliminary Results}
% Anything measured so far. If no results yet, say what the first test will be.
% [PROMPT: "describe any preliminary work done"] repo, tracks, interface doc, plans.

\todo{results or first planned test}

% \begin{figure}[h]\centering
% \includegraphics[width=0.8\textwidth]{figures/FILENAME.png}
% \caption{TODO}\label{fig:TODO}
% \end{figure}

%=====================================================================
\section{Evaluation Plan}
% Metrics, tracks (incl. unseen), number of seeds, how the comparison is kept fair.
% [PROMPT: "planned ... experiments and expected results"] -- the prompt asks for
% EXPECTED RESULTS too: add 2-3 sentences after the table (ppo_plan.tex §1 success).

\begin{center}
\begin{tabular}{lll}
\toprule
Metric & How measured & Why it matters\\
\midrule
\todo{} & & \\
\todo{} & & \\
\todo{} & & \\
\bottomrule
\end{tabular}
\end{center}

%=====================================================================
\section{Team Roles and Contributions So Far}
\begin{center}
\begin{tabular}{lll}
\toprule
Member & Responsibility & Done so far\\
\midrule
\todo{} & & \\
\todo{} & & \\
\todo{} & & \\
\todo{} & & \\
\bottomrule
\end{tabular}
\end{center}

%=====================================================================
\section{Timeline and Next Steps}
% Milestones until the 12-03 pre-report and 12-10 presentation.
% [PROMPT: "planned work"]

\todo{timeline}

%=====================================================================
\section{Risks and Open Questions}
% What could block you, and the fallback for each.
% Include the 2-3 biggest controller failure risks, one line each (risk -> fallback).
% FINAL REPORT (12-17): add a full "Failure Modes" section here, table with
% columns Failure | Detection | Response | Evidence, filled with real results.

\todo{risks}

\end{document}
